\documentclass{article}
\usepackage{amsmath,amssymb}
\usepackage[vocab={voting}]{leanify}
\newtheorem{lemma}{Lemma}
\begin{document}

\section{Notation}

We write $h(V,S)$ for the Hare quota for $V$ votes and $S$ seats.
We write $q(V,S)$ for the Droop quota for $V$ votes and $S$ seats.
We write $m(V)$ for the majority quota for $V$ votes.
We write $d(V,s)$ for the D'Hondt quotient of $V$ votes after $s$ seats.
We write $w(V,s)$ for the Webster quotient of $V$ votes after $s$ seats.
We write $c(n)$ for the number of pairwise contests among $n$ candidates.
We write $B(n,k)$ for the largest Borda score with $n$ candidates and $k$ voters.

\section{Quotas}

\begin{lemma}\label{quotas}
$h(100,4)=25$, $q(100,4)=21$ and $m(100)=51$.
\end{lemma}

\begin{lemma}\label{droop}
Let $V$ and $S$ be integers with $V\ge0$ and $S\ge1$. Then $(S+1)q(V,S)>V$.
\end{lemma}

\section{Divisors and counts}

\begin{lemma}\label{divisors}
$d(600,2)=200$, $w(600,2)=120$, $c(5)=10$ and $B(5,10)=40$.
\end{lemma}

\end{document}
